\documentclass[12pt,letterpaper]{article}



%%
%% Required packages:
%%
\usepackage{etex}
\usepackage{amsmath}
\usepackage{amssymb}
\usepackage{amstext}
\usepackage{amsthm}
\usepackage{fancyhdr,fancybox}
\usepackage{mathrsfs} % need for \mathscr command
\usepackage{multicol}
\usepackage[normalem]{ulem}
%\usepackage{pictex,pictexplus}
% \usepackage{pictexwd}
\usepackage{rotating}
\usepackage{url}
\usepackage{xfrac}
\usepackage{booktabs}
\usepackage{color,soul}
\usepackage{comment}

\usepackage{hyperref}
\hypersetup{
    colorlinks=true,     % false: boxed links; true: colored links
    linkcolor=blue,      % color of internal links
    citecolor=blue,      % color of links to bibliography
    filecolor=blue,      % color of file links
    urlcolor=blue        % color of external links
}


\usepackage{tikz}
\usetikzlibrary{backgrounds,calc}

%%
%% Common formatting commands
%%
\setlength{\parindent}{0in}
\setlength{\textwidth}{7in}
\setlength{\evensidemargin}{-0.25in}
\setlength{\oddsidemargin}{-0.25in}
\setlength{\parskip}{.5\baselineskip}
\setlength{\topmargin}{-0.5in}
\setlength{\textheight}{9in}

%               Problem and Part
\newcounter{problemnumber}
\newcounter{partnumber}
\newcounter{subpartnumber}
\newcommand{\Problem}{%
  \refstepcounter{problemnumber}%
  \setcounter{partnumber}{0}%
  \setcounter{subpartnumber}{0}%
  \item[\makebox{\hfil\textbf{\theproblemnumber.}\hfil}]%
  }
\newcommand{\InProblem}[2][1.6]{%
  \refstepcounter{problemnumber}%
  \setcounter{partnumber}{0}%
  \setcounter{subpartnumber}{0}%
  \textbf{\theproblemnumber.}\ \ \parbox[t]{#1in}{#2}%
}
\newcommand{\InSmallProblem}[1]{\InProblem[1.05]{#1}}
\newcommand{\NoProblem}[1]{%
  \mbox{\hfil\textbf{#1}\hfil}%
  }
\newcommand{\UnProblem}{%
  \refstepcounter{problemnumber}%
  \setcounter{partnumber}{0}%
  \setcounter{subpartnumber}{0}%
  \mbox{\hfil\textbf{\theproblemnumber.}\hfil}%
  }
\newcommand{\InFlexProblem}[1]{%
  \refstepcounter{problemnumber}%
  \setcounter{partnumber}{0}%
  \setcounter{subpartnumber}{0}%
  \textbf{\theproblemnumber.}\ \ #1%
}
%%  Part:
\newcommand{\Part}{%
  \renewcommand{\thepartnumber}{(\alph{partnumber})}%
  \refstepcounter{partnumber}
  \setcounter{subpartnumber}{0}%
  \item[(\alph{partnumber})]%
}
\newcommand{\InPart}[2][1.75]{%
  \renewcommand{\thepartnumber}{(\alph{partnumber})}%
  \refstepcounter{partnumber}%
  \setcounter{subpartnumber}{0}%
  (\alph{partnumber})\ \ \parbox[t]{#1in}{#2}%
}
\newcommand{\UnPart}{%
  \refstepcounter{partnumber}%
  \renewcommand{\thepartnumber}{(\alph{partnumber})}%
  \setcounter{subpartnumber}{0}%
  (\alph{partnumber})%
}
\newcommand{\InLargePart}[1]{\InPart[3.05]{#1}}
\newcommand{\InFullPart}[1]{\InPart[6.2]{#1}}
\newcommand{\InSmallPart}[1]{\InPart[1.05]{#1}}
%% SubPart:
\newcommand{\SubPart}{%
  \refstepcounter{subpartnumber}%
  \item[(\roman{subpartnumber})]%
  }
\newcommand{\InSubPart}[2][2.25]{%
  \stepcounter{subpartnumber}%
  (\roman{subpartnumber})\ \ \parbox[t]{#1in}{#2}%
}
\newcommand{\InSmallSubPart}[1]{\InSubPart[1.5]{#1}}
\newcommand{\InTinySubPart}[1]{%
  \stepcounter{subpartnumber}%
  (\roman{subpartnumber})\ \ \makebox{#1}%
}
\newcommand{\InMySubPart}[2][2.25]{%
  \stepcounter{subpartnumber}%
  \parbox[t]{#1in}{\makebox[0.3in][l]{(\roman{subpartnumber})}\ \ {#2}}%
}
\newcommand{\TrueFalse}{\textbf{T}\ \ \textbf{F}\ \ }
\newcommand{\TFPart}[1]{% 
  \stepcounter{partnumber}%
  \setcounter{subpartnumber}{0}%
  \item[(\alph{partnumber})] \TrueFalse\parbox[t]{5.5in}{#1}}

\newcommand{\Points}[1]{(#1 points) \ }
\newcommand{\Point}[1]{(#1 point) \ }


%% For Answers:
\newcounter{answernumber}
\newcommand{\Answer}{%
  \refstepcounter{answernumber}%
  \setcounter{partnumber}{0}%
  \setcounter{subpartnumber}{0}%
  \item[\makebox{\hfil\textbf{\theanswernumber.}\hfil}]%
  }


% Setting up the headers for the first page
%\chead{} % will default to what is typed in the file.
\fancyhead[L]{\textsc{Math 22a}}
\fancyhead[R]{\textsc{Fall 2024}}
\fancyfoot[R]{
	\vspace*{\fill}\footnotesize\textcopyright{} \emph{Harvard Math 22a}	
}
\renewcommand{\headrulewidth}{0pt}

%\fancyhead[C]{}
%	\chead{}	
%	\lhead{}
%	\rhead{}
%	\cfoot{}


% Putting copyright on the footer of each page
%\fancypagestyle{secondpage}{
%	\fancyhead[C]{TESTing again} % change this to be blank
%		%\chead{}	
%	\fancyhead[L]{bears} % change this to be blank
%	\fancyhead[R]{dogs} % change this to be blank
%	\fancyfoot[R]{ %keeps the footer the same
%		\vspace*{\fill}\footnotesize\textcopyright{} \emph{Harvard Math 22a}	
%	}
%}
%\renewcommand{\headrulewidth}{0pt}
%\pagestyle{fancy}
\pagestyle{empty}


%\lhead{}
%\rhead{}
%\rfoot{\vspace*{\fill}\footnotesize\textcopyright{} \emph{Harvard Math 22a}}
%\renewcommand{\headrulewidth}{0pt}
%\pagestyle{fancy}




%%
%% Common shortcut commands:
%%
\newcommand{\ds}{\displaystyle}
\newcommand{\sgn}{\operatorname{sgn}}
\renewcommand{\Pr}{\operatorname{P}}
\newcommand{\CPr}[3][{}]{\Pr_{\text{#1}}\left(#2\,|\,#3\right)}

\newcommand{\C}{\mathbb{C}}
\newcommand{\R}{\mathbb{R}}
\newcommand{\Z}{\mathbb{Z}}
\newcommand{\N}{\mathbb{N}}
\newcommand{\Q}{\mathbb{Q}}
\newcommand{\D}{\mathbb{D}}

\newcommand{\rref}{\operatorname{rref}}
\newcommand{\im}{\operatorname{im}}
\newcommand{\rank}{\operatorname{rank}}
\newcommand{\diag}{\operatorname{diag}}
\newcommand{\Span}{\operatorname{span}}
%\renewcommand{\vec}[1]{\mathbf{#1}}
\newcommand{\charpoly}{\mathscr{P}}
\newcommand{\nicefrac}[2]{\sfrac{#1}{#2}} % using xfrac package
\newcommand{\topology}{\mathcal{T}}
\newcommand{\Inn}{\operatorname{Inn}}

\newcommand{\blankbox}[2]{\fbox{\rule{#1}{0in}\rule{0in}{#2}}}

\newenvironment{myindentpar}[1]%
 {\begin{list}{}%
         {\setlength{\leftmargin}{#1}
          \setlength{\rightmargin}{\leftmargin}}%
         \item[]%
 }
 {\end{list}}
\newtheorem{theorem}{Theorem}
\newtheorem*{NStheorem}{Theorem}
\newtheorem{cor}[theorem]{Corollary}
\newtheorem*{claim}{Claim}
\newtheorem{defn}{Definition}

\newenvironment{amatrix}[1]{%
  \left(\begin{array}{@{}*{#1}{c}|c@{}}
}{%
  \end{array}\right)
}

%--------grstep
% For denoting a Gauss' reduction step.
% Use as: \grstep{\rho_1+\rho_3} or \grstep[2\rho_5 \\ 3\rho_6]{\rho_1+\rho_3}
\newcommand{\grstep}[2][\relax]{%
   \ensuremath{\mathrel{
       {\mathop{\longrightarrow}\limits^{#2\mathstrut}_{
                                     \begin{subarray}{l} #1 \end{subarray}}}}}}
\newcommand{\swap}{\leftrightarrow}


\usetikzlibrary{arrows}
\usepackage{wrapfig}

\makeatletter
\renewcommand*\env@matrix[1][*\c@MaxMatrixCols c]{%
	\hskip -\arraycolsep
	\let\@ifnextchar\new@ifnextchar
	\array{#1}}
\makeatother

\def\env@matrix{\hskip -\arraycolsep
	\let\@ifnextchar\new@ifnextchar
	\array{*\c@MaxMatrixCols c}}

\chead{\Large\textbf{Eigenvalues and Eigenvectors}}% 

\begin{document}
\thispagestyle{fancy}

\begin{center}
{\bf Review Questions}
\end{center}

\begin{enumerate}

\Problem {True or False Questions}

\begin{enumerate}
\item  If $A$ and $B$ are $n$ by $n$ matrices with $\det(A)=2$ and $\det(B)=3$, then $\det(A+B)=5$.
\item $k \text{det}(A) = \text{det}(kA)$
%\item $AB$ is invertible if and only if $A$ and $B$ are invertible.
\item If $A$ is row-equivalent to $B$, then \text{det}(A)=\text{det}(B).
\item If $A$ is invertible and $AB=0$, then $B=0$.
%\item If $T^2(x)=x$, then $T(x)=x$.
%\item $S$ and $T$ linear transformations.  If $S \circ T$ is one-to-one, then T is one-to-one.  Must $S$ be one-to-one?
%\item $S$ and $T$ linear transformations.  If $S \circ T$ is onto, then S onto.  Must $T$ be onto?
\item If $A^2=I$, then $A=I$ or $A=-I$.
\item If $A^2=0$, then $A=0$
%\item there exists $u,w,u \in \mathbb{R}^3$ linearly dependent, so that none of the 3 is a multiple of another.
\item $\text{det}(A) = \text{det}(B)$ implies $A-B$ is not invertible.
%\item Let $M$ be an $n$ by $n$ matrix with $n$ odd.  If $M^{T}=-M$, then $M$ is not invertible.
%\item If there exists a matrix $C$ so that $A= C^{-1} B C$, then \text{det}(A)=\text{det}(B).
%\item Let $T : \mathbb{R}^3 \to \mathbb{R}^3$ be rotation around $z$-axis, and $A$ be the standard matrix for $T$, then $\det{A}=1$.
%\item Let $T : \mathbb{R}^3 \to \mathbb{R}^3$ be projection onto the $yz$-plane, and $A$ be the standard matrix for $T$, then $\det{A}=0$.
%\item If $A$ is invertible and $A^{T}=-A$, then A is an $n$ by $n$ matrix where $n$ is even.
\item Let $T(x) =Ax$ and $S(x)=Bx$ and $A \sim B$, then $\text{im}(T) = \text{im}(S)$.
\item Suppose $A$ and $B$ have the same RREF, then the span of the columns of $A$ equals the span of the columns of $B$.
%\item If $Ax=b$ has a unique solution, then $A$ is a square matrix.
%\item If $A^2+4A+3I=0$, and $A$ is 3 by 3, then $A$ is invertible.
%\item If $A$ is $n$ by $n$ so that $A^2-A=0$, then $A=I$ or $A=0$.
%\item There is a 2 by 2 matrix $A$ with real entries so that $A^2 = -I$.
%\item There exists invertible 2 by 2 matrices $A$ and $B$ so that $\det(A+B)=\det(A)+\det(B)$.
%\item  For any 2 by 2 matrix A, there exists $\lambda \in \mathbb{R}$ so that $A+ \lambda I$ is singular.
\item  There is a linear transformation $T : \mathbb{R}^2 \to \mathbb{R}^2$ so that $\{ x : T(x)=0 \} = \{ T(x) : x \in \mathbb{R}^2 \}.$
\item   If $u, v, w$ are pairwise linearly independent, then $u, v, w$ are linearly independent. 
\end{enumerate}

\Problem 
\begin{enumerate}	
\item		Let $T\colon \R^3 \rightarrow \R^3$ be the transformation that scales all vectors by a factor of 3.
Let $A$ be the standard matrix for $T$.  What is $\det(A)$?
%\begin{solution}
%	We have that $T(\vec e_{1})=(3,0,0)$, $T(\vec e_{2})=(0,3,0)$ and $T(\vec e_{3})=(0,0,3)$.
%	So the standard matrix for $T$ is $A=\begin{bmatrix}3&0&0\\0&3&0\\0&0&3\end{bmatrix}$.
%	Since $A$ is triangular, the determinant of $A$ is the product of the diagonal entries: $\det(A)=27$.
%\end{solution}
%\blank
%\part[3]
\item Let $T\colon \R^3 \rightarrow \R^3$ be the transformation that reflects across the plane $x=z$.
Let $A$ be the standard matrix for $T$.  What is $\det(A)$?
%\begin{solution}
%	We have that $T(\vec e_{1})=\vec e_3, T(\vec e_{2})=\vec e_{2}$ and $T(\vec e_{3})=\vec e_{1}$.
%	So the standard matrix for $T$ is
%	$A=\begin{bmatrix}0&0&1\\0&1&0\\1&0&0\end{bmatrix}.$ Since $A$ is obtained from the identity
%	matrix by exchanging the first and third columns, $\det(A)=-1$.
%\end{solution}
%\blank
%\part[3]
\item Let $T\colon \R^3 \rightarrow \R^3$ be the transformation that projects to the plane $x=z$.
Let $A$ be the standard matrix for $T$. What is $\det(A)$?
%\begin{solution}
%	Since $T$ is not onto $\R^{3}$, $T$ is not invertible. So the standard matrix $A$ for $T$ is not
%	invertible. Thus, $\det(A)=0$.
%\end{solution}
%\blank
\end{enumerate}

\Problem Let \begin{equation*}
A=\begin{bmatrix}
1 & 0 & 1\\
3 & 3 & 4\\
2 & 2 & 3\\
\end{bmatrix}.
\end{equation*}  Find $A^{-1}$.

\Problem 
\begin{enumerate}
	\item Prove that if $T$ is one-to-one and $\vec{u_1}, \cdots, \vec{u_p}\in \mathbb{R}^n$ are linearly independent, then $T(\vec{u_1}), \cdots, T(\vec{u_p})$ are linearly independent in $\mathbb{R}^m$.
	\item In part (a), is the hypothesis that $T$ is one-to-one necessary?  Why or why not?
\end{enumerate}

\Problem Compute the determinant of the following matrices:
\begin{enumerate}
	\item
	\begin{equation*}
	A_1=\begin{bmatrix}
	1 & 1\\
	a_1 & 1\\
	\end{bmatrix}.
	\end{equation*}
	\item 
	\begin{equation*}
	A_2=\begin{bmatrix}
	1 & 1 & 1\\
	a_1 & 1 & 1\\
	a_2 & a_1 & 1\\
	\end{bmatrix}.
	\end{equation*}
	\item Find the determinant of the following matrix.  Use induction to prove your formula is correct.
	\begin{equation*}
	A_n=\begin{bmatrix}
	1 & 1 & \cdots & 1 & 1\\
	a_1 & 1 & \cdots & 1& 1\\
	\vdots &  & &&\vdots\\
	a_{n-1} & a_{n-2} &  \ddots &  1    & 1\\
	a_{n}     & a_{n-1} & \cdots & a_1 & 1\\
	\end{bmatrix}.
	\end{equation*}
\end{enumerate}
\Problem 
\begin{enumerate}
\item Prove $|\mathbb{N}|=|\{0\} \cup \mathbb{N}|$.
\item Prove $|\mathbb{N}|=|\{-2,-1\} \cup \mathbb{N}|$.
\item Prove $|[a,b]|=|[c,d]|$.
\end{enumerate}

\Problem Let  $T \colon \mathbb{P}_3 \rightarrow M_{2\times 2}$ be given by
			$$ T(a_0 + a_1 t + a_2 t^2 + a_3 t^3) = \begin{bmatrix} a_0-a_1 & a_2 - a_3 \\ a_1-a_0 & a_3 - a_2 \end{bmatrix}. $$ Find bases for the kernel and range of $L$.  Justify your answer.  Find the standard matrix of $L$ relative to the standard bases.
\Problem Let $U_{2 \times 2}$ be the space of upper triangular 2 by 2 matrices. Let $T : \mathbb{P}_2 \to U_{2 \times 2}$ be a linear transformation given by $$T(f) = \begin{bmatrix} f(0) & f'(1)\\  0 & f''(2)\\ \end{bmatrix}.$$   Find the matrix of $T$ relative to the standard bases.  Find the eigenvalues and eigenspaces of the matrix of $T$ relative to the standard bases.


\Problem Let $T$ be the linear transformation $\mathbb R^3 \to \mathbb R^3$ with matrix
	\[
	A = \begin{bmatrix}
	0 &
	1 &
	0 \\
	
	0 &
	1 &
	0 \\
	
	-1 &
	1 &
	1
	\end{bmatrix}.
	\]
	
	\begin{enumerate}
		\item Compute the dimension of the kernel $T$ and a basis for the kernel.
		
		\item Compute the rank of $A$ and a basis for the image of $T$.
	\end{enumerate}
	
%\Problem
%	Let
%	$$
%	A = \begin{bmatrix}-10 &-18\\9&17\end{bmatrix}.
%	$$
%	\begin{enumerate}
%		\item Compute the eigenvalues of $A$.
%		\item Compute a basis for each eigenspace of $A$.
%		\item Compute a matrix $B$ such that $B^3=A$.
%	\end{enumerate}
	
\Problem Let $T:\mathbb{P}_2\to\mathbb{R}^3$ be defined by $$T(p(t)) = \begin{bmatrix}p(1)\\p(2)\\p(3)\end{bmatrix}$$ for every polynomial $p(t) = a_2t^2 + a_1t + a_0$ in $\mathbb{P}_2$.
	
	\begin{enumerate}
		\item Show that $T$ is a linear transformation.
		\item Find a matrix $M$ such that $T(p(t)) = M[p(t)]_{\mathcal{B}}$ where $\mathcal{B} = \{1,t,t^2\}$.
		\item Show that for all $c_1, c_2, c_3 \in \mathbb{R}$ there is a polynomial $p(t) \in \mathbb{P}_2$ such that $p(1) = c_1$, $p(2) = c_2$, and $p(3) = c_3$.
	\end{enumerate}
	

	

%\Problem
%	\begin{enumerate}
%		\item
%		Give an example of a real $n \times n$ matrix $A$ with eigenvalues $\frac{-1+i\sqrt{3}}{2}$ and $\frac{-1-i\sqrt{3}}{2}$.
%		\item
%		Give a complete list of all diagonalizable matrices with characteristic polynomial $(\lambda - 2)^4$. You must prove your list is complete.
%		\item
%		Find the algebraic multiplicities, and the dimension of the eigenspaces, of all the eigenvalues of the linear transformation $D:\mathbb{P}_2\to\mathbb{P}_2$ which sends each polynomial $p(t) = a_2t^2 + a_1t+ a_0$ in $\mathbb{P}_2$ to its derivative $p'(t) = 2a_2t+a_1$.
%		
%	\end{enumerate}
	
%\Problem
%	Let $\mathcal{B} = \{\mathbf{b}_1,\mathbf{b}_2,\mathbf{b}_3\}$ be a basis for $\mathbb{R}^3$. Suppose that $A$ and $B$ are $3 \times 3$ matrices such that $\mathbf{b}_1,\mathbf{b}_2,\mathbf{b}_3$ are eigenvectors of both $A$ and $B$. Show that $AB = BA$.
	
\Problem
	Let $A$ be an $n\times n$ matrix which is not invertible, and define $K_r$ to be the nullspace of $A^r$.  That is, $K_1 = \mbox{Nul}(A)$, $K_2 = \mbox{Nul}(A^2)$, etc.
	
	\begin{enumerate}
		\item Prove that, for each $j=1,2,3,\ldots$, $K_j$ is a subspace of $K_{j+1}$, so that 
		$$
		K_1\subseteq K_2 \subseteq K_3 \subseteq K_4 \subseteq\cdots.
		$$
		\item Prove that, if for some $r$ we have $K_r = K_{r+1}$, then $K_r=K_n$ for all $n\geq r$.
	\end{enumerate}
	\Problem More True or False Questions.
	\begin{enumerate}
		 \item There is a $2 \times 4$ matrix $A$ with rank $1$ and $$\mathrm{Nul}(A) = \mathrm{Span}\left\{\begin{bmatrix}1\\0\\0\\1\end{bmatrix},\begin{bmatrix}1\\0\\1\\0\end{bmatrix},\begin{bmatrix}1\\1\\0\\0\end{bmatrix},\begin{bmatrix}3\\1\\1\\1\end{bmatrix}\right\}.$$
		\item Let $V$ and $W$ be vector spaces and $T : V \to W$ be a one-to-one linear transformation.  If $\{ T(v_1), \cdots, T(v_k)\}$ is a basis for $W$, then $\{ v_1, \cdots, v_k \}$ is a basis for $V$.
		\item The set $\{A\in M_{n \times n}(\mathbb{R}) : A \text{ is invertible. } \}$ is a subspace of $M_{n \times n} (\mathbb{R})$.
		\item Let $A$ be an $n$ by $n$ matrix such that the sum of each row is equal to 1.  Then $A-I$ is not invertible.
%		\item Let $E_1$ and $E_2$ be eigenspaces of distinct eigenvlaues of a matrix $A$.  Then $E_1 \cap E_2 = \{ 0 \}$.
		\item  Let $W_1$ and $W_2$ be 3-dimensional subspaces of $\mathbb{R}^5$, then $W_1$ and $W_2$ have a non-zero vector in common.
		\item $V$, $W$ subspaces of $\mathbb{R}^n$, then $V \cup W$ is a subspace.
		\item  Let $W_1$ and $W_2$ be subspaces of $\mathbb{R}^n$ such that $\dim(W_1)\leq \dim(W_2)$.  Then $W_1\subset W_2$.
	%	\item If $x$ is an eigenvector of a matrix $A$, then $x$ is also an eigenvector of $A^3$.
		\item If $A$ is an $n \times n$ matrix and $x, y$ are two eigenvectors of $A$, then $x + y$ is also an eigenvector of $A$.
		\item If $V$ and $W$ are vector spaces of dimension $n$ and $T : V \to W$ is a one-to-one linear transformation, then $T$ is onto.
%		\item If $A$ and $B$ both have the characteristic polynomial $-\lambda^3+7\lambda^2-10\lambda$, then $A$ and $B$ are similar matrices.
	\end{enumerate}


\end{enumerate}

\newpage
\chead{\Large\textbf{Eigenvalues and Eigenvectors -- Answers / Solutions}}%
\lhead{}
\rhead{}
\thispagestyle{fancy}

\begin{enumerate}

  \Answer 
  
  \begin{enumerate}
  \item {\bf False.} Let $A=\begin{bmatrix} 2 & 0\\ 0 & 1\\ \end{bmatrix}$ and $B=\begin{bmatrix} 1 & 0\\ 0 & 3\\ \end{bmatrix}$. Then $\det{(A+B)} = \det \begin{bmatrix} 3  & 0\\ 0 & 4\\ \end{bmatrix}\neq 5$.
  
  \item {\bf False.} Let $k=2$ and $A=I_2$, then $k\cdot \det{A}=2$ and $\det{(kA)}=4$.
  
  \item {\bf False.} Let $A=I_2$ and $B=\begin{bmatrix} 0 & 1\\ 1 & 0\end{bmatrix}$, then $\det{A}=1$ and $\det{B}=-1$ and $A$ is still row equivalent to $B$.
  
  \item {\bf True.} Since $A$ is invertible, we can multiply both sides be $A^{-1}$ on the left to get $B=A^{-1} (0)=0$.
  
  \item {\bf False.} Let $A=\begin{bmatrix} 0 & 1\\ 1 & 0\\ \end{bmatrix}$, then $A^2=I_2$ and $A\neq I_2$ and $A\neq -I_2$.
  
  \item {\bf False.} Let $A=\begin{bmatrix} 0 & 1\\ 0 & 0\\ \end{bmatrix}$, then $A^2=0$ and $A\neq 0$.
  
  \item {\bf False.} Let $A=\begin{bmatrix} 1 & 0\\ 0 & 0\\ \end{bmatrix}$ and $B=\begin{bmatrix} 0 & 0\\ 0 &-1\\ \end{bmatrix}$, then $A-B=\begin{bmatrix} 1 &0 \\ 0 & 1\\ \end{bmatrix}$. Thus $\det{A}=\det{B}$ and $A-B$ is invertible.
  
  \item {\bf False.} Let $A=\begin{bmatrix} 0 & 1\\ 0& 0\\ \end{bmatrix}$ and $B=\begin{bmatrix} 0 & 0 \\ 0 & 1\\\end{bmatrix}$, then $A\sim B$. Notice that $\text{im}(T)=\left\{\begin{bmatrix} x\\ 0 \end{bmatrix} : x\in \mathbb{R} \right\}$ and $\text{im}(S) =\left \{ \begin{bmatrix} 0 \\ y\\ \end{bmatrix} : y \in \mathbb{R} \right \}$. Thus their images are not the same.
  
  \item {\bf False.} Let $A=\begin{bmatrix} 0 & 1\\ 0& 0\\ \end{bmatrix}$ and $B=\begin{bmatrix} 0 & 0 \\ 0 & 1\\\end{bmatrix}$, then $A$ and $B$ have the same reduced row echelon form, but the span of the columns is different.
  
  \item {\bf True.} Let $T: \mathbb{R}^2 \to \mathbb{R}^2$ be given by $T(\vec{x})= \begin{bmatrix} 0 & 1\\ 0 &0 \\ \end{bmatrix}$. Then $\{ x: T(x)=0\} = \text{Span}\left( \begin{bmatrix} 1\\ 0\\ \end{bmatrix} \right)$ and 
  $\{T(x): x\in \mathbb{R}^2\} =\text{im}(T)=\text{Span}\left( \begin{bmatrix} 1\\ 0\\ \end{bmatrix} \right)$.
  
  \item {\bf False.} Let $\vec{u} =\begin{bmatrix} 1\\ 0 \end{bmatrix}$, $\vec{v} =\begin{bmatrix} 0\\ 1 \end{bmatrix}$, and $\vec{w} =\begin{bmatrix} 1\\ 1 \end{bmatrix}$. Then each of the three pairs of vectors will be linearly independent, but all three vectors together are linearly dependent.
  
  \end{enumerate}
  
  \Answer 
  
  \begin{enumerate}
  
  \item 		We have that $T(\vec e_{1})=(3,0,0)$, $T(\vec e_{2})=(0,3,0)$ and $T(\vec e_{3})=(0,0,3)$.
			So the standard matrix for $T$ is $A=\begin{bmatrix}3&0&0\\0&3&0\\0&0&3\end{bmatrix}$.
			Since $A$ is triangular, the determinant of $A$ is the product of the diagonal entries: $\det(A)=27$.

  
  \item 		We have that $T(\vec e_{1})=\vec e_3, T(\vec e_{2})=\vec e_{2}$ and $T(\vec e_{3})=\vec e_{1}$.
			So the standard matrix for $T$ is
			$A=\begin{bmatrix}0&0&1\\0&1&0\\1&0&0\end{bmatrix}.$ Since $A$ is obtained from the identity
			matrix by exchanging the first and third columns, $\det(A)=-1$.
  
  \item  Since $T$ is not onto $\R^{3}$, $T$ is not invertible. So the standard matrix $A$ for $T$ is not
			invertible. Thus, $\det(A)=0$.
  
  \end{enumerate}
  
  
  \Answer We augment $A$ by $I_3$ and row-reduce to find $A^{-1}$. Thus

\begin{align*}
	\begin{pmatrix}[ccc|ccc]
		1 & 0 & 1 & 1 & 0 & 0\\
		3 &3 & 4 & 0 & 1 & 0\\
		2 &2 & 3 & 0 & 0 &1\\
	\end{pmatrix}& \sim  \begin{pmatrix}[ccc|ccc]
		1 & 0 & 1 & 1 & 0 & 0\\
		0 &3 & 1 & -3 & 1 & 0\\
		0 & 2 & 1 & -2 & 0 &1\\
	\end{pmatrix}\\ &\sim \begin{pmatrix}[ccc|ccc]
		1 & 0 & 1 & 1 & 0 & 0\\
		0 &3 & 1 & -3 & 1 & 0\\
		0 &0 & 1 & 0 & -2 &3\\
	\end{pmatrix}\\ &\sim
\begin{pmatrix}[ccc|ccc]
	1 & 0 & 0 & 1 & 2 &-3 \\
	0 &1 & 0 & -1 & 1 & -1\\
	0 &0 & 1 & 0 & -2 & 3\\
\end{pmatrix},
\end{align*}
and hence $A^{-1}=\begin{bmatrix}
1 & 2 & -3\\
-1 & 1 & -1\\
0 & -2 & 3\\
\end{bmatrix}.$
  
  \Answer 
  
    \begin{enumerate}
  
  \item Consider the vector equation
\begin{equation}
c_1 T(\vec{u_1})+ \cdots + c_n T(\vec{u_n}) = 0.
\end{equation}
Since $T$ is linear, we have 
\begin{equation}
 T(c_1 \vec{u_1}+ \cdots + c_n \vec{u_n}) = 0.
\end{equation}
Since $T$ is one-to-one, we know that if $T(\vec{x})=0$, then $\vec{x}=0$.

Thus we have 
\begin{equation}
 c_1 \vec{u_1}+ \cdots + c_n \vec{u_n} = 0.
\end{equation}
Since $\vec{u_1}, \cdots, \vec{u_n}$ are linearly independent, we must have $c_1=c_2=\cdots=c_n=0$.

Therefore $T(\vec{u_1}), \cdots, T(\vec{u_n})$ are linearly independent.
  
  \item The hypothesis that $T$ is one-to-one is necessary.  Suppose $T$ is the zero transformation.  Let $\vec{u_1}, \cdots, \vec{u_n}\in \mathbb{R}^n$ be linearly independent, then $T(\vec{u_j})=0$ for all $j=1,\cdots, n$.  And hence $T(\vec{u_1}), \cdots, T(\vec{u_n})$ are linearly dependent.  Therefore $T$ is a counterexample to the claim without the injectivity hypothesis.
  
  \end{enumerate}
  
  \Answer 
  
    \begin{enumerate}
  
  \item In the first case, we get $\det{A_1}=1-a_1$.
  
  \item In the second case, we get $\det{A_2}=1-2a_1-a_1^2=(1-a_1)^2$.
  
  \item From the first few cases, it seems like $\det{A_n}=(1-a_1)^n$.
  
  \begin{proof}
  We prove the result by induction.
  
  {\bf Base case:} We proved the base case in part (a).
  
  {\bf Inductive Step:} We assume $\det{A_n} = (1-a_1)^n$, and we want to prove $\det{A_{n+1}}=(1-a_1)^{n+1}$.
  
  We expand the determinant about the first column to get
  $$\det{A_{n+1}}= \det (A_{n+1})_{11} - a_1 \det (A_{n+1})_{21}+\dots+(-1)^{n+1+1}a_n \det (A_{n+1})_{1(n+1)}.$$
  
  Now notice $(A_{n+1})_{11}=A_n$ and $(A_{n+1})_{21}=A_n$. For $(A_{n+1})_{j1}$, $j>2$, we have a repeated row of ones, and hence $\det{(A_{n+1})_{j1}}=0$ when $j>2$. Using these facts and the induction hypothesis, we get 
    $$\det{A_{n+1}}=\det A_n -a_1 \det  A_n= (1-a_1)^n-a_1(1-a_1)^n=(1-a_1)^{n+1}.$$
    
    Thus the result holds by mathematical induction.
  \end{proof}
  
  \end{enumerate}
  
  \Answer
  
    \begin{enumerate}
  
  \item We want a bijection from $\mathbb{N}$ to $\{ 0\} \cup \mathbb{N}$. Define $f: \mathbb{N} \to \{0\}\cup \mathbb{N}$ by $f(n)=n-1$. Note $f$ is a bijection.
  
  \item Define a bijection from $\mathbb{N}$ to $\{-2, -1\}\cup \mathbb{N}$ by $g(k)=\begin{cases}
  -2 & \text{if } k=1\\
    -1 & \text{if } k=2\\
     k-2 & \text{if } k>2\\
  \end{cases}$.

 \item  Define $h: [a,b] \to [c,d]$ by $h(x)= c+ \frac{d-c}{b-a}(x-a)$. 
  \end{enumerate}

\Answer 			A vector $p(t) = a_0 + a_1 t + a_2 t^2 + a_3 t^3$ will be in the kernel if $a_0 - a_1 = a_2 - a_3 = 0$.
				Thus a basis for the kernel will be 
					$\left\{ 1+t, t^2+t^3 \right\}$.
					
					Let $\mathcal{B} = \{ 1,t,t^2,t^3\}$ be a basis for $\mathbb{P}_3$, and $\mathcal{C} = \left\{ \begin{bmatrix} 1 & 0 \\ 0 & 0 \end{bmatrix},\begin{bmatrix} 0 & 1 \\ 0 & 0 \end{bmatrix}, \begin{bmatrix} 0 & 0 \\ 1 & 0 \end{bmatrix}, \begin{bmatrix} 0 & 0 \\ 0 & 1 \end{bmatrix} \right\}$
			be a basis for $M_{2 \times 2}$.
			
							We need to find the matrix
					$\left[ \begin{array}{ccc} [T(b_1)]_{\mathcal{C}} & \cdots & [T(b_n)]_{\mathcal{C}} \end{array} \right]$:				\begin{align*}
					T(b_1) &= T(1) = \begin{bmatrix} 1 & 0 \\ -1 & 0 \end{bmatrix} \\
					T(b_2) &= T(t) = \begin{bmatrix} -1 & 0 \\ 1 & 0 \end{bmatrix} \\
					T(b_3) &= T(t^2) = \begin{bmatrix} 0 & 1 \\ 0 & -1 \end{bmatrix} \\
					T(b_4) &= T(t^3) = \begin{bmatrix} 0 & -1 \\ 0 & 1 \end{bmatrix}
				\end{align*}
				Thus the matrix is
				\[
					\begin{bmatrix} 1 & -1 & 0 & 0 \\ 0 & 0 & 1 & -1 \\ -1 & 1 & 0 & 0 \\ 0 & 0 & -1 & 1 \end{bmatrix}.
				\]
				
		From the previous part, we know that the range of $T$ is the span of 
				\[
					\left\{ \begin{bmatrix} 1 & 0 \\ -1 & 0 \end{bmatrix}, \begin{bmatrix} -1 & 0 \\ 1 & 0 \end{bmatrix}, \begin{bmatrix} 0 & 1 \\ 0 & -1 \end{bmatrix}, \begin{bmatrix} 0 & -1 \\ 0 & 1 \end{bmatrix} \right\}.
				\]

				Removing the dependent entries, we find a basis
				\[
					\mathcal{D} = \left\{ \begin{bmatrix} 1 & 0 \\ -1 & 0 \end{bmatrix}, \begin{bmatrix} 0 & 1 \\ 0 & -1 \end{bmatrix} \right\}
				\]
					
					
\Answer Evaluating $T$ on the standard basis for $\mathbb{P}_2$ yields: 
$$T\left(1 \right) = \begin{bmatrix} 1 & 0\\ 0 & 0\end{bmatrix};  \; \;
T\left( t \right) = \begin{bmatrix} 0 & 1\\ 0 & 0\end{bmatrix}; \; \;
T\left( t^2 \right) = \begin{bmatrix} 0 & 2\\ 0 & 2\end{bmatrix}.$$

Let $\mathcal{E}$ denote the standard basis of $U_{2 \times 2}$, then 
$$[T\left(1 \right)]_{\mathcal{E}} = \begin{bmatrix} 1\\ 0\\  0\\\end{bmatrix};  \; \;
[T\left( t \right)]_{\mathcal{E}} = \begin{bmatrix} 0 \\ 1\\ 0\\\end{bmatrix}; \; \;
[T\left( t^2 \right)]_{\mathcal{E}} = \begin{bmatrix} 0 \\ 2 \\ 2\end{bmatrix}.$$

Thus the matrix of $T$ relative to the standard coordinates is given by
$$M=\begin{bmatrix}
1 & 0 & 0\\
0 & 1 & 2\\
0 & 0 & 2\\
\end{bmatrix}.$$

We proved in class that the eigenvalues of an upper triangular matrix are given by the diagonal entries, and hence the eigenvalues of $T$ are $1$ and $2$.

We first find the eigenspace corresponding to the 1-eigenvalue.  Thus, we have
$$E_1 = N(M- I)= \text{Span}\left\{  \begin{bmatrix} 1\\0\\0\\\end{bmatrix}, \begin{bmatrix} 0\\1\\0\\\end{bmatrix}\right\}.$$

Similarly, we find the eigenspace for $\lambda=2$.  That is, we have
$$E_2 = N(M-2I)= \text{Span}\left\{  \begin{bmatrix} 0\\2\\1\\ \end{bmatrix}\right\}.$$

Since we have 3 linearly independent eigenvectors, the diagonalization theorem implies $T$ will act diagonally in the bases coming from the eigenvectors.  Namely, we have the eigenbasis $$\left\{  \begin{bmatrix} 1\\0\\0\\\end{bmatrix}, \begin{bmatrix} 0\\1\\0\\\end{bmatrix}, \begin{bmatrix} 0\\2\\1\\ \end{bmatrix}\right\}.$$  

This gives the corresponding bases $\mathcal{P} =\left\{1, t, 2t+t^2 \right \}$ in $\mathbb{P}_2$ and $ \mathcal{B}= \left\{  \begin{bmatrix} 1 & 0\\ 0 & 0\end{bmatrix}, \begin{bmatrix} 0 & 1\\ 0 & 0\end{bmatrix}, \begin{bmatrix} 0 & 2\\ 0 & 1\end{bmatrix}\right\}$ for $U_{2 \times 2}$.  Thus the matrix $M$ of $T$ relative to $\mathcal{P}$ and $\mathcal{B}$ is given by
$$M=\begin{bmatrix}
1 & 0 & 0\\
0 & 1 & 0\\
0 & 0 & 2\\
\end{bmatrix}.$$
\Answer 
\begin{enumerate}
\item The kernel of $T$ is the nullspace of $A$, so we row reduce $A$ to
find 
$$\begin{bmatrix} 1 & 0 & -1 \\ 0 & 1 & 0 \\ 0 & 0 & 0\end{bmatrix}.$$
Hence a basis for the kernel consists of the one vector
$$\begin{bmatrix}1\\0\\1\end{bmatrix},$$ 
which means that the kernel of $T$ is one-dimensional.
\item We already row reduced $A$ and we found two pivots, which means that
the rank of $A$ is $2$.

The image of $T$ is the column space of $A$ and a basis for that space
consists of the first two columns of $A$:
$$\begin{bmatrix}0\\0\\-1\end{bmatrix}, \begin{bmatrix}1\\1\\1\end{bmatrix}.$$
\end{enumerate}

\Answer
\begin{enumerate}
\item We must show that $T((p+q)(t)) = T(p(t)) + T(q(t))$ and $T((cp)(t)) =
cT(p(t))$ hold for all polynomials $p(t)$ and $q(t)$ in $\mathbb{P}_2$
and every scalar $c \in \mathbb{R}$.

Both equalities are immediate since $(p+q)(t)$ is defined to equal
$p(t) + q(t)$ and $(cp)(t)$ is defined to equal $c(p(t))$ for every
value of $t$.
\item Write $p(t) = a_2t^2 + a_1t + a_0$ so that $[p(t)]_{\mathcal{B}}
= \begin{bmatrix}a_0\\a_1\\a_2\end{bmatrix}$. 
Then $$\begin{bmatrix}p(1)\\p(2)\\p(3)\end{bmatrix}
= \begin{bmatrix}a_2 + a_1 + a_0\\a_24 + a_12 + a_0\\a_29 + a_13 +
  a_0\end{bmatrix}
= \begin{bmatrix}1&1&1\\1&2&4\\1&3&9\end{bmatrix}\begin{bmatrix}a_0\\a_1\\a_2\end{bmatrix},$$
which shows that $M = \begin{bmatrix}1&1&1\\1&2&4\\1&3&9\end{bmatrix}$ works.
\item We are asked to solve the equation $T(p(t))
= \begin{bmatrix}c_1\\c_2\\c_3\end{bmatrix}$ for the polynomial $p(t)
= a_2t^2 + a_1t + a_0$.
By part~(b), this equation is equivalent to the matrix
equation $$\begin{bmatrix}c_1\\c_2\\c_3\end{bmatrix}
= \begin{bmatrix}1&1&1\\1&2&4\\1&3&9\end{bmatrix}\begin{bmatrix}a_0\\a_1\\a_2\end{bmatrix}.$$
This equation always has a solution since the matrix $M$ is invertible.
\end{enumerate}


\Answer
\begin{enumerate}
\item We need to show that $K_j \subseteq K_{j+1}$.
So suppose $x \in K_j$, that is $A^jx = 0$. Then $$A^{j+1}x
=A(A^jx) = A\mathbf0 = \mathbf0,$$ which shows that $x \in K_{j+1}$.
\item Suppose that $K_r = K_{r+1}$. We prove the result by induction on $n \geq 1$ that
$K_r = K_{r+n}.$

The base case $n = 1$ is given to us!

For the induction step, suppose we have shown that $K_r = K_{r+n}$.
To show that $K_r = K_{r+n+1}$ we can show instead that $K_{r+n} =
K_{r+n+1}$.
We already know that $K_{r+n} \subseteq K_{r+n+1}$ so it suffices to
show that $K_{r+n+1} \subseteq K_{r+n}$.

Suppose that $x \in K_{r+n+1}$, that is $A^{r+n+1}x = \mathbf0$.
Then $A^{r+1}(A^n x) = \mathbf0$, which means that $A^n x \in
K_{r+1}$. Since $K_{r+1} = K_r$, we conclude that $A^n x \in K_r$ and
hence $A^r(A^n x) = \mathbf0$. Since $A^rA^n = A^{n+r}$, we conclude
that $x \in K_{n+r}$.
\end{enumerate}
\Answer
\begin{enumerate}
\item True. An example of such a matrix is $$A
= \begin{bmatrix}1&-1&-1&-1\\0&0&0&0\end{bmatrix}.$$
Indeed, a basis for the null space of this matrix consists of the
three
vectors $$\begin{bmatrix}1\\0\\0\\1\end{bmatrix},\begin{bmatrix}1\\0\\1\\0\end{bmatrix},\begin{bmatrix}1\\1\\0\\0\end{bmatrix},$$
which is also a basis of 
$$\mathrm{Span}\left\{\begin{bmatrix}1\\0\\0\\1\end{bmatrix},\begin{bmatrix}1\\0\\1\\0\end{bmatrix},\begin{bmatrix}1\\1\\0\\0\end{bmatrix},\begin{bmatrix}3\\1\\1\\1\end{bmatrix}\right\}$$
since the last spanning vector is redundant.
\item True.  We need to show that $\{ v_1, \cdots, v_k \}$ is linearly independent and spans $V$.  Let $v \in V$, then $T(v) \in W$.  Since $\{ T(v_1), \cdots, T(v_k)\}$ is a basis for $W$, there exists $c_1, \cdots, c_k \in \mathbb{R}$ so that $$T(v) = c_1 T(v_1) +\cdots T(v_k).$$  Since $T$ is injective and linear, we get $v = c_1 v_1+\cdots +c_k v_k$.  Thus, $v_1, \cdots, v_k$ span $V$.

We are left to check linear independence.  Consider the vector equation $d_1 v_1 + \cdots d_k v_k = 0_V$.  Applying $T$ to both sides and using the linearity of $T$ yields $$d_1 T(v_1) + \cdots d_k T(v_k) = 0.$$  Since $T(v_1), \cdots, T(v_k)$ are linearly independent, we know that $d_1=d_2 = \cdots = d_k =0$.  Thus $v_1, \cdots, v_k$ are linearly independent, and hence they form a basis. 
\item False.  The zero matrix is not invertible, so the set cannot be a subspace.
\item True.  Since the sum of each row of $A$ equals 1, we know that 1 is an eigenvalue of $A$ with eigenvector of all 1's.  Thus $A-I$ must have non-trivial null space, and hence by the invertible matrix theorem, $A-I_n$ is not invertible.
\item True.  Let $v \in E_1 \cap E_2$.  Let $E_j$ be the eigenspace corresponding to $\lambda_j$ for $j=1,2$.  Then $\lambda_1 v =\lambda_2 v$ since $\lambda_1\neq \lambda_2$, we must have $v=0$.
\item True. Let $x_1,x_2,x_3$ be a basis for $W_1$ and $y_1, y_2, y_3$ be a basis for $W_2$. Since $x_1, x_2, x_3, y_1, y_2, y_3$ is 6 vectors in $\mathbb{R}^5$, the vectors must be linearly dependent. Thus there are scalars, not all zero, such that $$c_1 x_1 +c_2 x_2+ c_3 +c_4 y_1+ c_5 y_2+c_6 y_3 =0,$$ and hence  $$c_1 x_1 +c_2 x_2 + c_3 x_3 = -c_4 y_1-c_5 y_2-c_6 y_3.$$ Thus $c_1 x_1 +c_2 x_2 + c_3 x_3$ is a non-zero element in $W_1 \cap W_2$.
\item False. Let $V=\text{Span}\left\{\begin{bmatrix} 1\\ 0\\ \end{bmatrix} \right\}$ and $W=\text{Span}\left\{\begin{bmatrix} 0\\ 1\\ \end{bmatrix} \right\}$. Now $\begin{bmatrix} 1\\ 0\\ \end{bmatrix}, \begin{bmatrix} 0\\ 1\\ \end{bmatrix} \in V\cup W$, but $\begin{bmatrix} 1\\ 0\\ \end{bmatrix}+ \begin{bmatrix} 0\\ 1\\ \end{bmatrix}=\begin{bmatrix} 1\\ 1\\ \end{bmatrix}\not \in V \cup W.$
\item False. Let $W_1=\text{Span}\left\{\begin{bmatrix} 1\\ 0\\ \end{bmatrix} \right\}$ and $W_2=\text{Span}\left\{\begin{bmatrix} 0\\ 1\\ \end{bmatrix} \right\}.$ Then $\dim W_1 \leq \dim W_2$, but $W_1$ is not a subset of $W_2$.
%\item True. Suppose that $Ax = \lambda x$. Then $$A^3x = A^2(Ax) =
%A^2(\lambda x) = \lambda A^2 x = \lambda A(A x) = \lambda A(\lambda x)
%= \lambda^2 A x = \lambda^3 x,$$ which shows that $x$ is an
%eigenvector of $A^3$ with eigenvalue $\lambda^3$.
%\item False. Take $A = \begin{bmatrix}1&0\\0&-1\end{bmatrix}$. Then $x
%= \begin{bmatrix}1\\0\end{bmatrix}$ and $y
%= \begin{bmatrix}0\\1\end{bmatrix}$ are both eigenvectors of $A$, but
%$x+y$ is not an eigenvector of
%$A$ since
%$$A(x+y)
%= \begin{bmatrix}1&0\\0&-1\end{bmatrix}\begin{bmatrix}1\\1\end{bmatrix}
%= \begin{bmatrix}1\\-1\end{bmatrix}$$ is not a multiple of $x + y$.
\item True. Because $T$ is one-to-one, the dimension of the image of $T$
must be $n$. Since $W$ has dimension $n$, the image of $T$ must equal $W$.
%\item True. The characteristic polynomial has three distinct roots: $0$,
%$2$, $5$. Therefore $A$ and $B$ are both similar to the diagonal
%matrix $\begin{bmatrix}0&0&0\\0&2&0\\0&0&5\end{bmatrix}$, which means
%that they are also similar to eachother.
\end{enumerate}
        
          

\end{enumerate}  

\end{document}


%%% Local ables: 
%%% mode: latex
%%% TeX-master: t
%%% End: 
